\section{Discussion, Limitations, and Conclusion}
\label{sec:discussion}

Instruction quality is the most stable explanatory component in this pilot.
The \(B-A\) comparison changes instructions while retaining the two-call,
single-persona structure, and its direction persists across scenario families
and trimming.  Claims about structured reasoning should therefore compare
against instruction-matched baselines rather than crediting an architectural
label for gains caused by task specification.

This point extends beyond prompt engineering as an implementation detail.
Instructions determine which assumptions must be checked, how missing evidence
is treated, and whether a system recognizes abstention as a scientifically
valid outcome.  In a validity-sensitive task, omissions in the task contract
can surface as methodological errors rather than merely stylistic differences.
Evaluations should therefore treat instruction content as an explicit
experimental factor.

Staging may still matter by preventing rare, high-cost errors.  C achieves
the best aggregate outcome, yet most of its incremental advantage over B
comes from one structural-break cell; trimming reverses the sign and the two
families disagree.  This supports a hypothesis about tail-error prevention,
not a broad average-case benefit from additional calls.  Failure-sensitive
evaluation is essential because near-zero medians can coexist with sharply
different means.

The pattern suggests a distinction between average-case refinement and
tail-risk control.  Additional stages may contribute little on cells where the
eligible action is obvious yet still matter when a review stage catches an
expensive invalid continuation.  The present data support that hypothesis only
as a mechanism for future study: one influential cell cannot establish a
general tail-protection effect, and the scenario families disagree on the
average \(C-B\) direction.

The \(D-C\) comparison is likewise bounded.  The tested specialized workflow
underperforms unified staging, but it changes persona scope, intermediate
interfaces, message flow, and context.  Specialized agents may remain useful
when roles possess genuinely different tools, information, objectives, or
verification duties.  Our result characterizes two workflow bundles and
cannot establish a general single-agent advantage.

Specialization may be valuable when roles possess genuinely different tools,
information, incentives, or verification duties
\citep{HongEtAl2024MetaGPT}.  Those conditions are largely absent here: every
stage addresses the same statistical task with the same model family and
underlying evidence environment.  Adding boundaries in this setting may add
coordination and interface costs without adding useful diversity.  A future
factorial design would need to vary personas, context partitioning, and
intermediate schemas separately before attributing the observed difference.

These findings also clarify what ERI contributes.  Typed evidence and
fail-closed routing make invalid decisions and abstentions measurable, but
neither is claimed as novel.  The empirical contribution is the ordering of
comparisons: instruction content is varied before call organization, and the
specialized workflow is evaluated only after an instruction-matched unified
staged condition.  This sequence narrows attribution without pretending to
recover independent causal effects for every workflow feature.

It also motivates reporting practices for structured scientific systems.
Aggregate accuracy alone cannot distinguish an unnecessary abstention from an
invalid inferential claim, and mean loss alone can conceal concentration.
Complete denominators, component outcomes, paired cell differences, and
prespecified sensitivity analyses reveal whether a workflow improves ordinary
cases, prevents rare severe errors, or simply changes its failure mode.

\paragraph{Design implications.}
The results suggest a practical ordering for scientific-AI development.
Begin with a precise task contract and an instruction-matched baseline before
adding calls or roles.  When additional stages are introduced, preserve
component-level outcomes so that a gain can be attributed to corrected
evidence handling, eligibility assessment, or final routing.  Add specialized
roles only when the task supplies genuine heterogeneity---different tools,
information, objectives, or verification responsibilities---rather than
assuming that persona labels alone create useful diversity.

Evaluation should also distinguish uniform improvement from prevention of
rare catastrophic errors.  Median regret is near zero for several conditions
even though their means differ sharply.  A system can look reliable on
ordinary cases while occasionally making an invalid high-loss decision.
Paired cell traces, scenario-specific summaries, leave-one-extreme
sensitivities, and explicit failure retention are therefore central to
interpreting the architecture, not optional robustness checks after a headline
average has been selected.

The scope is narrow: 18 cells, two synthetic scenario families, three seeds,
one hosted model and provider, and only two scenario clusters.  Hosted-model
outputs may vary even under nominally deterministic settings
\citep{AtilEtAl2025Nondeterminism}.  The loss profile is benchmark-specific,
the deterministic policy was co-developed with the benchmark, and the
conditions are not a complete factorial design.  Broader claims require
replication across models, domains, and less formalizable evidence; additional
limitations and provenance qualifications are in
Supplement Section K.  In particular, the two-cluster intervals are not
conventional uncertainty estimates for a population of tasks, and the
benchmark-specific loss weights affect the magnitude of routing regret.

The four conditions are also not a complete factorial decomposition.  B and C
match substantive instructions and nominal output capacity, but differ in
call count, context accumulation, and persisted intermediate state.  C and D
differ in persona specialization together with message partitioning,
intermediate schemas, and role boundaries.  The study therefore cannot assign
an independent causal coefficient to staging or persona identity.  It reports
the cleanest instruction contrast and qualified implementation contrasts for
the remaining bundles.

The benchmark itself is narrow and highly structured.  It focuses on
Sharpe-ratio inference under synthetic dependence, heteroskedasticity,
distributional, stability, and provenance conditions.  Results may differ
where evidence is unstructured, ontologies are ambiguous, tools are
heterogeneous, expertise is genuinely distributed, or adversarial interaction
is part of the task.  The tested specialized workflow's result must not be
read as evidence that multi-agent systems are generally inferior.

Two additional operational qualifications matter.  All live conditions use
one hosted model and provider configuration, and hosted services can vary
under nominally deterministic settings
\citep{AtilEtAl2025Nondeterminism}.  Two units terminated at provider-schema
validation and are retained using the preregistered fail-closed convention.
The recovered evidence package preserves primary outcomes, aggregates,
authorization, freeze records, and transcripts, but not a complete original
request/response and token ledger; operational quantities without artifact
support are therefore excluded from the canonical scientific claims.

The deterministic boundary should be read with the same restraint.  Its
approximately \(0.0013\) mean regret demonstrates what direct encoding can do
in a task whose evidence ontology and admissibility rules were built alongside
the benchmark.  It does not show that deterministic policies can be specified
reliably in open-ended domains.  LLMs may remain useful when evidence is
unstructured or the decision ontology is incomplete, even if a formalized
final boundary is executed deterministically.

The deterministic result also depends on co-development.  Its evidence
ontology, diagnostic acquisition, priorities, and abstention rules encode
substantial knowledge aligned with the benchmark.  It is thus an engineered
reference for what complete formalization can achieve, not evidence that
similarly reliable rules can be specified in advance for open-ended domains.
The study does not evaluate autonomous policy synthesis.

Finally, routing regret is a policy-weighted benchmark quantity.  The
registered loss makes invalid selection and harmful non-abstention expensive
because those outcomes threaten scientific validity.  Different applications
may place different costs on delay, abstention, interval width, or invalid
claims and can therefore change numerical rankings.  Reporting categorical
validity outcomes alongside regret makes this dependence visible but does not
eliminate it.

Within these limits, the pilot supports three descriptive findings.  Richer
instructions produce the clearest and most stable improvement.  Unified
four-call execution has the strongest aggregate profile, but its advantage over
instruction-matched two-call execution is concentrated in one severe
structural-break error and does not persist uniformly across sensitivities.
The tested specialized workflow does not improve on unified staging, although
bundled architectural differences prevent attributing that outcome solely to
role specialization.

The immediate next test is replication rather than architectural escalation:
additional scenario families, independently constructed admissibility
policies, multiple models and providers, and tasks with less structured
evidence.  Such studies can ask whether the stable instruction effect
persists, whether staging repeatedly prevents tail failures, and when genuine
role heterogeneity offsets interface costs.  Until then, every conclusion
remains benchmark-specific and pilot-descriptive.

The broader contribution is methodological.  Workflow comparisons should match
instruction content, retain failures in the denominator, and report paired
cell-level outcomes and concentration analyses before attributing gains to
additional calls or specialized agents.  \ERI{} demonstrates this principle in a
scientific setting where invalid method selection and harmful non-abstention
can be distinguished from ordinary answer error.  A workflow that produces more
answers is not necessarily safer, and a system that abstains is not
necessarily less capable.  By separating evidence acquisition, admissibility,
action, and failure accounting, the benchmark makes those distinctions
observable and auditable.
